%% kapa: exam=1 round=20 year=2009 period=2 form=A subject=appraisal_law subject_ko=부동산관계법규 seq=13 number=53
%% source: Q-Net 감정평가사 기출문제 (cst00309) artlSeq=5024178 file=extracted/1008/5024178/2009년 제20회 감정평가사 1차시험 기출문제/2교시.hwp
%% conversion: kordoc 4.24.51 -> tools/gen1.py (auto); review=proofread-sonnet; flags=none
\begin{question}{53}
국토의 계획 및 이용에 관한 법령상 공공시설에 관한 설명으로 옳지 \uline{않은} 것은?
\begin{choices}
  \item 도로ㆍ공원ㆍ철도ㆍ수도ㆍ행정청이 설치하는 주차장은 공공시설에 해당한다.
  \item 행정청이 아닌 자가 개발행위허가를 받아 새로 공공시설을 설치한 경우, 용도폐지된 종래의 공공시설은 개발행위허가를 받은 자에게 무상으로 귀속된다.
  \item 개발행위허가를 받은 자가 행정청이 아닌 경우, 개발행위허가를 받은 자가 새로 설치한 공공시설은 그 시설을 관리할 관리청에 무상으로 귀속된다.
  \item 행정청이 도시계획시설사업에 의하여 새로 공공시설을 설치한 경우, 새로 설치한 공공시설은 그 시설을 관리할 관리청에 무상으로 귀속된다.
  \item 개발행위허가를 받은 자가 행정청인 경우, 행정청에 귀속된 공공시설의 처분으로 인한 수익금을 도시계획사업 외의 목적에 사용하여서는 아니 된다.
\end{choices}
\end{question}
\answer{2}
